\chapter{Security hardening: connecting the trust boundaries}

Earlier chapters taught authentication, object-level authorization, CSRF and CORS protection, typed SQL, AppFs, shared runtime state, and outbound egress policy separately. The detailed production configuration key reference comes later; here we summarize the security boundaries that configuration must enforce. In production handoff, these are not independent features but one security model.

Security hardening therefore does not introduce new syntax. Its purpose is to verify that every external value and every sensitive operation crosses a clear trust boundary and that no shortcut bypasses compiler, runtime, or server policy.

\section{The five most important trust boundaries}

A typical RWLang application should treat at least these boundaries separately:

\begin{longtable}{L{0.25\textwidth}L{0.31\textwidth}L{0.34\textwidth}}
\toprule
Boundary & Untrusted side & Authority / validation \\
\midrule
HTTP input & path, query, header, form, JSON & typed route/input, parser limits, CSRF/CORS \\
Identity & cookie, login fields, client role & trusted auth backend, session, MFA \\
Data & client value, record identifier & typed query, bind, transaction, authorization \\
Filesystem & filename, MIME, upload bytes & AppFs confinement, server-owned destination, limits \\
External network & URL, DNS, peer, response & named egress target, CIDR/port/TLS/timeout policy \\
\bottomrule
\end{longtable}

The most important mental rule is: \textbf{data does not become trusted merely because it already passed through an earlier layer}. An authenticated user's form data is still user input; a valid record identifier does not imply authorization; a host resolved by DNS does not prove that the actual peer IP is allowed.

\section{HTTP edge: Host, proxy, and parser}

With direct HTTPS, the RWLang server itself is the public HTTP trust boundary. Behind Apache or Nginx there are two layers: the public proxy and the RWLang listener.

In production, forwarding headers may be accepted only from explicitly trusted proxy networks. A directly client-supplied \code{X-Forwarded-For}, \code{X-Forwarded-Proto}, or similar header must not become authority.

The public host should also come from trusted configuration. Do not construct redirect targets from an arbitrary client \code{Host} header.

The RWLang HTTP parser is intentionally strict and rejects ambiguous request forms. Examples include duplicate critical headers, obs-fold, and bodies on GET/HEAD. Request \code{Transfer-Encoding} and unknown connection tokens are also rejected. The developer rule is simple: \textbf{do not attempt to loosen edge-parser invariants from application code}.

\subsection*{Fail-closed decisions and security audit}

A trust-boundary denial should not exist only as a client response. The server uses the same server-generated request ID for access logs and, where a security decision occurred, a structured \code{security_audit} event. Important classifications are:

\begin{longtable}{L{0.30\textwidth}L{0.12\textwidth}L{0.46\textwidth}}
\toprule
Boundary & HTTP & Audit category/outcome \\
\midrule
forwarding metadata & 400 & \code{proxy/invalid_forwarding} \\
Host pinning & 421 & \code{host/mismatch} \\
HTTPS policy & 426 & \code{transport/https_required} \\
missing auth on API & 401 & \code{auth/unauthorized} \\
CSRF & 403 & \code{csrf/denied} \\
CORS & 403 & \code{cors/denied} \\
Origin/Fetch Metadata & 403 & \code{origin/denied} \\
route/object policy & 403 & \code{policy/forbidden} \\
rate limit & 429 & \code{rate_limit/denied} \\
\bottomrule
\end{longtable}

Do not put tokens, cookies, credentials, full URLs, or request bodies into the security audit. Category/outcome says \emph{which trust boundary decided}; use the request ID to correlate access/server logs for detailed request forensics.

\section{TLS and secret handling}

Private keys, database credentials, Redis URLs, LDAP credentials, and other secrets must not live in:

\begin{itemize}
  \item source code;
  \item the Git repository;
  \item command-line arguments;
  \item plaintext TOML secret values;
  \item access/server/audit logs.
\end{itemize}

The production pattern is a \code{*_file} reference:

\begin{lstlisting}
[tls]
cert_file = "/run/secrets/tls/fullchain.pem"
key_file = "/run/secrets/tls/privkey.pem"

[database]
url_file = "/run/secrets/rwlang/database-url"
\end{lstlisting}

Configuration states \emph{where} the secret is; the secret value itself remains behind a separate permission boundary. Secret files should be readable by the service user and as inaccessible as practical to everyone else.

\section{Identity: authentication is not authorization}

Successful login proves only that a trusted principal exists. It does not prove that the principal may access a particular article, invoice, or admin operation.

The correct layers are:

\begin{enumerate}
  \item authentication: who is the user;
  \item route authorization: may the user enter this handler;
  \item object authorization: may the user access this particular record;
  \item business invariant: is this operation allowed in the current state.
\end{enumerate}

Client-supplied role, owner, tenant, or price is never authority. Roles come from the trusted auth backend, owner comes from a trusted record, and business decisions come from server-side state.

\Needspace{0.38\textheight}
A typical object-level mutation order is:

\begin{lstlisting}[language=RWLang]
let article = Article.byId(db, id)?;
authorize article owner authorUsername or role Publisher;

transaction db {
    Article.publishChanged(tx, id, version)?;
    audit Article id action publish;
        from article.status
        to ArticleStatus.Published
}
\end{lstlisting}

The authorization guard does not replace optimistic locking, and optimistic locking does not replace authorization.

\section{Session, CSRF, and CORS together}

A session cookie is not CSRF protection by itself. State-changing form or JSON requests require session-bound CSRF protection; JSON POST may carry the token in a dedicated header.

Treat browser-side protection as defense in depth:

\begin{lstlisting}
session-bound CSRF token
+ same-origin Origin/Referer check
+ Sec-Fetch-Site cross-site deny
\end{lstlisting}

Open CORS only to an explicit exact-origin allowlist. Credentialed CORS must never use a wildcard. CORS is not authentication or authorization; it only controls which browser origins can read a response.

\section{SQL and business mutation}

Compiler-owned SQL plus typed bind parameters are among V1's strongest security foundations. User input must not enter SQL string concatenation, dynamic column names, or dynamic table names.

Mutations run inside transactions. A transaction is not only a consistency tool: it also matters for security because multiple writes performed after authorization should succeed or fail together, and a business audit should belong to the same state transition.

The runtime DB credential should follow least privilege. Migration credentials are a separate operator capability; the web server should not run continuously with schema-modification privileges.

\section{Uploads, media, and AppFs}

Every byte of an uploaded file is untrusted. Client filename and MIME type are metadata only; do not use them as storage keys, paths, or file-type authority.

The production baseline is:

\begin{itemize}
  \item server-owned destination directory;
  \item byte limit;
  \item for images, pixel/decode limit and byte-based format validation;
  \item AppFs confinement;
  \item no symlink/magic-link escape;
  \item runtime data separated from release/static files;
  \item only the required \code{r}/\code{w}/\code{c} capability.
\end{itemize}

Do not use the same directory for public static root and user uploads merely because both need to be served over HTTP. Their lifecycle, trust, and cache policies are different.

\section{Redis and shared runtime state}

Redis is runtime-owned infrastructure, not business source of truth. Sessions, TOTP replay state, rate limits, and public cache may live there; business records should not move there simply because Redis is ``faster.''

With multiple instances, missing shared session/rate-limit state can produce different security decisions depending on which node receives a request. Shared-state policy is therefore part of the production topology.

When Redis fails, security-relevant functions should not become automatically permissive. Fail closed for rate limits, replay protection, and session validation.

\section{Outbound egress and SSRF}

The external network is a separate trust boundary. V1 has no unrestricted application-side \code{HttpClient}; the trusted integration adapter uses named targets.

Target policy constrains host, address range, port, TLS, timeouts, and response size together. After DNS resolution, every A/AAAA address is checked, then the actual peer address is validated again after connection.

Forbidden pattern:

\begin{lstlisting}
user input URL -> generic server-side fetch
\end{lstlisting}

Recommended pattern:

\begin{lstlisting}
user/business choice
        |
        v
named integration target
        |
        v
trusted host/CIDR/port/TLS policy
\end{lstlisting}

The user may choose a business option, but not an arbitrary network destination.

In V1, egress is a capability of the trusted Rust integration layer, not a general \code{.rw} HTTP API. Therefore a fail-closed \code{IntegrationError::Policy} denial is not automatically a webserver \code{security_audit} event: the adapter/operator integration must log the target name and minimal outcome if forensic evidence is required. Never log secrets, credentials, arbitrary full URLs, or response bodies there either.

\section{Resource limits are security controls too}

Security is not only confidentiality and authorization. A web application is also vulnerable if one request can consume unbounded memory, CPU, DB time, or upload bandwidth.

The defensive layers work together:

\begin{itemize}
  \item request/parser hard limits;
  \item upload and response size limits;
  \item query/result limits and timeouts;
  \item rate limits;
  \item runtime resource profiles;
  \item process/systemd/cgroup hard limits.
\end{itemize}

Do not request a larger resource profile ``just in case.'' Larger budgets should be a production policy decision and an auditable exception.

\section{Log hygiene and audit integrity}

Logs are themselves a data-exfiltration surface. Do not log credentials, cookies, CSRF/TOTP secrets, or entire form/JSON bodies.

Security events need minimal structured fields: request ID, route, actor when appropriate, decision type, and status. Business audit should be transactional and tied to the application state change.

A request ID is a correlation identifier, not authority. The server should generate its own; do not trust a client-supplied identifier.

\section{Process and filesystem hardening}

Run the production service as an unprivileged user. Release directories and configuration should be runtime read-only; only genuinely writable areas should be writable.

Recommended separation:

\begin{lstlisting}
/usr/local/bin/rwlang-server
/usr/local/bin/rwlang-cli
/usr/local/etc/rwlang/       # trusted config, read-only
/srv/rwlang/releases/        # immutable releases
/srv/rwlang/data/            # declared persistent app data
/run/secrets/rwlang/         # secret files
/var/log/rwlang/             # logs
\end{lstlisting}

Systemd hardening, cgroup limits, and RWLang resource policy are not alternatives; they are separate defensive layers.

\section{Dependency and build supply chain}

The deployed RWLang server is a binary, so the dependency graph is release state too. A committed \code{Cargo.lock} and \code{--locked} builds make the dependency surface more reproducible and reviewable.

After an intentional dependency update, do more than compile:

\begin{lstlisting}
./tools/refresh-lock.sh
./verify.sh
./tools/dependency-audit.sh
\end{lstlisting}

RustSec audit provides evidence about known advisories; it does not replace code review or upstream changelog review. Major/minor migrations of parser, TLS, crypto, auth, or database dependencies should not be hidden inside unrelated feature diffs.

\section{Production security gate}

Minimum checks before production handoff:

\begin{enumerate}
  \item HTTPS and public-host policy are verified.
  \item Reverse-proxy trust is open only to explicit CIDRs.
  \item Secrets enter the process only through file references.
  \item Runtime DB credentials are least privilege; migration credentials are separate.
  \item Session/auth backend matches the single-node or multi-instance topology.
  \item Every state-changing browser operation has CSRF protection.
  \item CORS uses exact allowlists; no credentialed wildcard exists.
  \item Object authorization uses trusted principal/role and stable owner authority.
  \item SQL is compiler-owned statement + typed binds only.
  \item Upload uses a server-owned destination and restricted AppFs capability.
  \item No business source of truth would disappear if Redis were empty after restore.
  \item Outbound integrations reach the network only through named egress policy.
  \item Request, upload, DB, runtime, and process resource limits are configured.
  \item Logs contain no secrets or full sensitive request payloads.
  \item Security-negative tests, config check, and deployment smoke tests are green.
\end{enumerate}

\section{A code-review thinking pattern}

Five questions catch many security problems early in every new feature:

\begin{enumerate}
  \item \textbf{What comes from outside?} -- Which field, header, file, record identifier, or external response is untrusted?
  \item \textbf{Where does it become typed data?} -- Route parser, form/JSON schema, model, or trusted adapter?
  \item \textbf{Who is the authority?} -- Session backend, database record, trusted config, or egress policy?
  \item \textbf{What happens on failure?} -- Fail closed, rollback, cleanup, 4xx/5xx, audit?
  \item \textbf{What limits the resource?} -- Size, time, count, rate, or process budget?
\end{enumerate}

If these questions do not have short, unambiguous answers, the feature's security model is probably not finished.
